\section{Evaluation Methodology}
\label{sec:method}

\subsection{Single-Margin Simulator}

Two robots travel on perpendicular paths through one crossing, and robot A may
brake and wait for robot B.  The monitored margin is a time headway,
$M=\lvert\hat g\rvert-k_\sigma\sigma-G_{\mathrm{safe}}$, where $\hat g$ is the
estimated difference between the arrival times of B and A at the crossing,
$\sigma$ is a range-dependent uncertainty allowance, and
$G_{\mathrm{safe}}=0.5$\,s is the required headway; additive Gaussian noise
perturbs each measurement.  An optimistic hidden bias $s f^{p}$ delays B's
estimated arrival, where $s$ is the fault severity in seconds, $f\in[0,1]$ is
B's remaining range fraction, and the exponent $p$ sets how the bias decays:
$p<1$ reveals the conflict late, $p=1$ linearly, and $p>1$ early.  Recovery is
rate-limited braking with capability $a_{\max}$.  The predictive rule estimates
the margin slope with the Theil--Sen estimator \cite{sen1968theilsen} over the
eight most recent samples.  To avoid dividing by a fitted slope near zero, it
projects only when the fitted decay exceeds $10^{-3}$ margin units per second
($10^{-6}$ in the coordination simulator) and otherwise uses its $M\le0$
fallback.

Two model-specific quantities bound the latest useful trigger; neither is a
general safety certificate.  A continuous-time braking calculation gives the
latest trigger that stops A outside the wider crossing zone, $8.65$\,s with the
default parameters; it is conservative relative to the simulator's collision
condition.  Deterministic replay on the simulator's $0.1$\,s control grid finds
$8.8$\,s as the latest tested forced-trigger time that avoids a collision in the
default conflict episode; this boundary is empirical and grid-dependent.

\subsection{Three-Robot Coordination Simulator}

Three robots approach one shared zone at a nominal speed of $1.2$\,units/s; the
zone is one unit long, so the nominal occupancy is $0.83$\,s.  In the conflict
condition, start distances of $(12,12.5,13)$ units put adjacent nominal entry
times only $0.42$\,s apart, so untreated trajectories overlap physically.  In the
clear condition, start distances of $(12,14,16)$ units separate entries by
$1.67$\,s.  A collision is recorded when interpolated zone-occupancy intervals
overlap, and a separate headway violation when two entry times are less than
$5/6$\,s apart.  Each robot's estimated arrival time includes a decaying hidden
bias with per-robot multipliers $(0,1,2)$ and the exponent family used above.
For each pair of robots, the monitored margin is the absolute difference between
their estimated arrival times, minus the required headway and both uncertainty
allowances, plus measurement noise.  The fixed rule triggers when any pairwise
margin is at most $\rho$.  The predictive rule keeps a separate Theil--Sen slope
history for each pair and triggers when a margin is nonpositive or its projected
time to zero is at most $T_h$.

Both rules invoke the same recovery, an explicit one-shot serialization state
machine.  Robots start in the \textsc{approaching} state.  At the first trigger,
any robot already \textsc{crossing} is placed first, the others are ordered by
their current estimated arrival times, and the order is committed for the
episode.  Every later robot enters \textsc{waiting} with all earlier robots as
blockers and returns to \textsc{approaching} only after every blocker is
\textsc{cleared}.  Pairwise slope histories are reset at commitment, and no robot
yields twice.  The design isolates the effect of trigger timing on a fixed
recovery; it is not a general multi-robot planner.  \Cref{tab:parameters} lists
the main parameters of both simulators.

\begin{table}[!t]
  \centering
  \caption{Simulator and Evaluation Parameters}
  \label{tab:parameters}
  \footnotesize
  \setlength{\tabcolsep}{4pt}
  \begin{tabular}{@{}lll@{}}
    \toprule
    Parameter & Single margin & Coordination \\
    \midrule
    Robots & 2 & 3 \\
    Time step & $0.1$\,s & $0.1$\,s \\
    Nominal speed & $1.2$\,units/s & $1.2$\,units/s \\
    Required headway & $0.5$\,s & $5/6$\,s \\
    Measurement-noise s.d. & $0.05$\,s & $0.05$\,s \\
    Slope window & 8 samples & 8 samples \\
    $a_{\max}$ (units/s$^2$) & $0.5,1.0,2.5$ & $0.5,1.0,2.5$ \\
    Degradation exponent $p$ & $0.7,1.0,1.5$ & $0.7,1.0,1.5$ \\
    Fault severity $s$ (s) & $2,4,6,8,10$ & $2,6,10$ \\
    Conflict scenarios & 45 & 27 \\
    Pad and horizon grid & $0$--$4$, step $0.25$ & $0$--$4$, step $0.25$ \\
    Tuning / test replicates & 25 / 150 & 20 / 60 \\
    Held-out conflict episodes & 6,750 per arm & 1,620 per arm \\
    \bottomrule
  \end{tabular}
\end{table}


\subsection{Evaluation Protocol}

\emph{Exploratory sweeps.}  The dose-response experiment compares an
uncertainty-blind arrival-gap trigger, the current-margin trigger at $M=0$, and
the predictive rule with $T_h=2.0$\,s, using 40 paired noise seeds at each of
seven severities.  The recovery-capability and degradation-shape sweeps use 25
paired seeds at each of five severities, or 125 episodes per plotted setting and
rule; the shape sweep fixes $a_{\max}=0.8$.  For each setting, the fixed-pad
candidates are $\rho\in\{0,0.5,1.0,1.5,2.0,2.5,3.0,4.0\}$\,s, and the displayed
comparator is the candidate with the lowest mean successful elapsed time among
those for which every episode succeeds on the same setting and seed grid.
Selection and reporting therefore use the same simulations, so these figures are
mechanism checks rather than out-of-sample estimates.

\emph{Seed-disjoint single-margin evaluation.}  The evaluation crosses three
recovery capabilities, three degradation exponents, and five severities, giving
the 45 conflict scenarios of \cref{tab:parameters}.  The fixed pad and the
predictive horizon are both searched from $0$ to $4$ in steps of $0.25$.  Each
candidate receives 25 tuning replicates per scenario (1,125 episodes).  Among
candidates with no tuning collision, selection minimizes the mean conflict
elapsed time and breaks ties toward the smaller parameter.  Testing then uses
150 new replicates per scenario, or 6,750 episodes per conflict arm.  Clear-path
behavior is evaluated at $a_{\max}=1.0$ over the 15 exponent--severity
combinations, for 2,250 episodes per arm; a yield in this condition is a false
yield.  The original $T_h=2.0$\,s rule is retained as a diagnostic arm, but the
primary comparison uses the separately tuned horizon.

\emph{Seed-disjoint coordination evaluation.}  Tuning crosses the same three
capabilities and exponents with three severities, giving 27 scenarios, and
searches both parameters over the same grid.  Each candidate receives 20 tuning
replicates per scenario (540 episodes).  A candidate qualifies only if every
tuning episode reaches the exit without collision, headway violation, timeout,
or serialization breach; selection minimizes the mean successful makespan and
breaks ties toward the smaller parameter.  Each selected rule and a
no-intervention reference then receive 60 new replicates per scenario and
condition, or 1,620 episodes.

\emph{Seeds.}  Every split--scenario--replicate tuple maps to a distinct
deterministic seed, giving 10,125 single-margin and 3,780 coordination seeds
with no overlap between tuning and testing.  A tuple's seed is shared across
candidates and rules to preserve pairing.  The split therefore holds out noise
realizations only; the simulators and scenario grids remain fixed.

\subsection{Statistical Reporting}

Failure counts are reported before any timing comparison.  For each exact
scenario, we compute a two-sided 95\% Wilson score interval for the collision
probability \cite{wilson1927}.  We report the largest upper endpoint across the
pointwise intervals; this maximum is neither a simultaneous confidence bound nor
an estimate for a population of scenarios, and pooled counts summarize a fixed
factorial design rather than one homogeneous binomial experiment.  Elapsed time
and makespan are summarized only over successful episodes.  A paired percentile
bootstrap with 20,000 resamples, clustered by test replicate
\cite{efron1993bootstrap}, is computed only when every compared pair succeeds.
If failure outcomes differ, we report descriptive successful-only means and
suppress the inferential timing comparison.

\subsection{Connection to the Timing Model}

The recovery-capability sweep varies $a_{\max}$ as a proxy for recovery
duration, and the exponent sweep varies how quickly the margin degrades.  The
simulations do not measure the episode-level distributions of $\tau_r$,
$\delta$, $L$, and $\Delta$, and they introduce no observation delay.  They test
behavior consistent with the timing argument but do not validate
\cref{prop:boundary,prop:homogeneity} by measuring all of their variables.
